%-----------------------------------------------------------------------------------------------------------------------------------------------%
%	The MIT License (MIT)
%
%	Copyright (c) 2015 Jan Küster
%
%	Permission is hereby granted, free of charge, to any person obtaining a copy
%	of this software and associated documentation files (the "Software"), to deal
%	in the Software without restriction, including without limitation the rights
%	to use, copy, modify, merge, publish, distribute, sublicense, and/or sell
%	copies of the Software, and to permit persons to whom the Software is
%	furnished to do so, subject to the following conditions:
%	
%	THE SOFTWARE IS PROVIDED "AS IS", WITHOUT WARRANTY OF ANY KIND, EXPRESS OR
%	IMPLIED, INCLUDING BUT NOT LIMITED TO THE WARRANTIES OF MERCHANTABILITY,
%	FITNESS FOR A PARTICULAR PURPOSE AND NONINFRINGEMENT. IN NO EVENT SHALL THE
%	AUTHORS OR COPYRIGHT HOLDERS BE LIABLE FOR ANY CLAIM, DAMAGES OR OTHER
%	LIABILITY, WHETHER IN AN ACTION OF CONTRACT, TORT OR OTHERWISE, ARISING FROM,
%	OUT OF OR IN CONNECTION WITH THE SOFTWARE OR THE USE OR OTHER DEALINGS IN
%	THE SOFTWARE.
%	
%
%-----------------------------------------------------------------------------------------------------------------------------------------------%


%============================================================================%
%
%	DOCUMENT DEFINITION
%
%============================================================================%

%we use article class because we want to fully customize the page and dont use a cv template
\documentclass[10pt,A4]{article}	


%----------------------------------------------------------------------------------------
%	ENCODING
%----------------------------------------------------------------------------------------

%we use utf8 since we want to build from any machine
\usepackage[utf8]{inputenc}		

\usepackage{hyperref}

\makeatletter
\hypersetup{
	%pagebackref=true,
	colorlinks=true,       		% false: boxed links; true: colored links
	linkcolor=black,          	% color of internal links
	citecolor=black,        		% color of links to bibliography
	filecolor=magenta,      		% color of file links
	urlcolor=black,
	bookmarksdepth=4
}
\makeatother
%----------------------------------------------------------------------------------------
%	LOGIC
%----------------------------------------------------------------------------------------

% provides \isempty test
\usepackage{xifthen}

%----------------------------------------------------------------------------------------
%	FONT
%----------------------------------------------------------------------------------------

% some tex-live fonts - choose your own

%\usepackage[defaultsans]{droidsans}
%\usepackage[default]{comfortaa}
%\usepackage{cmbright}
\usepackage[default]{raleway}
%\usepackage{fetamont}
%\usepackage[default]{gillius}
%\usepackage[light,math]{iwona}
%\usepackage[thin]{roboto} 

% set font default
\renewcommand*\familydefault{\sfdefault} 	
\usepackage[T1]{fontenc}

% more font size definitions
\usepackage{moresize}		


%----------------------------------------------------------------------------------------
%	PAGE LAYOUT  DEFINITIONS
%----------------------------------------------------------------------------------------

%debug page outer frames
%\usepackage{showframe}			


%define page styles using geometry
\usepackage[a4paper]{geometry}		

% for example, change the margins to 2 inches all round
\geometry{top=1cm, bottom=1cm, left=1.5cm, right=1.5cm} 	

%use customized header
\usepackage{fancyhdr}				
\pagestyle{fancy}

%less space between header and content
\setlength{\headheight}{-5pt}		


%customize entries left, center and right
%\lhead{}
%\chead{}
%\rhead{}


%indentation is zero
\setlength{\parindent}{0mm}

%----------------------------------------------------------------------------------------
%	TABLE /ARRAY DEFINITIONS
%---------------------------------------------------------------------------------------- 

%for layouting tables
\usepackage{multicol}			
\usepackage{multirow}

%extended aligning of tabular cells
\usepackage{array}

\newcolumntype{x}[1]{%
>{\raggedleft\hspace{0pt}}p{#1}}%


%----------------------------------------------------------------------------------------
%	GRAPHICS DEFINITIONS
%---------------------------------------------------------------------------------------- 

%for header image
\usepackage{graphicx}

%for floating figures
\usepackage{wrapfig}
\usepackage{float}
%\floatstyle{boxed} 
%\restylefloat{figure}

%for drawing graphics		
\usepackage{tikz}				
\usetikzlibrary{shapes, backgrounds,mindmap, trees}


%----------------------------------------------------------------------------------------
%	Color DEFINITIONS
%---------------------------------------------------------------------------------------- 

\usepackage{color}

%accent color
\definecolor{sectcol}{RGB}{0,150,255}

%dark background color
\definecolor{bgcol}{RGB}{110,110,110}

%light background / accent color
\definecolor{softcol}{RGB}{225,225,225}


%============================================================================%
%
%
%	DEFINITIONS
%
%
%============================================================================%

%----------------------------------------------------------------------------------------
% 	HEADER
%----------------------------------------------------------------------------------------

% remove top header line
\renewcommand{\headrulewidth}{0pt} 

%remove botttom header line
\renewcommand{\footrulewidth}{0pt}	  	

%remove pagenum
\renewcommand{\thepage}{}	

%remove section num		
\renewcommand{\thesection}{}			

%----------------------------------------------------------------------------------------
% 	ARROW GRAPHICS in Tikz
%----------------------------------------------------------------------------------------

% a six pointed arrow poiting to the left
\newcommand{\tzlarrow}{(0,0) -- (0.2,0) -- (0.3,0.2) -- (0.2,0.4) -- (0,0.4) -- (0.1,0.2) -- cycle;}	

% include the left arrow into a tikz picture
% param1: fill color
%
\newcommand{\larrow}[1]
{\begin{tikzpicture}[scale=0.58]
	 \filldraw[fill=#1!100,draw=#1!100!black]  \tzlarrow
 \end{tikzpicture}
}

% a six pointed arrow poiting to the right
\newcommand{\tzrarrow}{ (0,0.2) -- (0.1,0) -- (0.3,0) -- (0.2,0.2) -- (0.3,0.4) -- (0.1,0.4) -- cycle;}

% include the right arrow into a tikz picture
% param1: fill color
%
\newcommand{\rarrow}
{\begin{tikzpicture}[scale=0.7]
	\filldraw[fill=sectcol!100,draw=sectcol!100!black] \tzrarrow
 \end{tikzpicture}
}



%----------------------------------------------------------------------------------------
%	custom sections
%----------------------------------------------------------------------------------------

% create a coloured box with arrow and title as cv section headline
% param 1: section title
%

\newcommand{\eventspace}{
	\vspace{0.05cm}
}

\newcommand{\sectionspace}{
	\vspace{0.1cm}
}
\newcommand{\cvsection}[1]
{
	\begin{center}
		\large\textcolor{sectcol}{\textbf{#1}}
	\end{center}
	\sectionspace
}

%create a coloured arrow with title as cv meta section section
% param 1: meta section title
%
\newcommand{\metasection}[2]
{
%\begin{tabular*}{1\textwidth}{r r}
\footnotesize{#2} \hspace*{\fill} \footnotesize{#1}\\[1pt]
%\end{tabular*}
}

%----------------------------------------------------------------------------------------
%	 CV EVENT
%----------------------------------------------------------------------------------------

% creates a stretched box as cv entry headline followed by some paragraphs about 
% the work you did
% param 1:	event time i.e. 2014 or 2011-2014 etc.
% param 2:	event name (what did you do?)
% param 3:	institution (where did you work / study)
% param 4:	list of paragraphs
%
\newcommand{\cvevent}[4]
{

\begin{tabular*}{1\textwidth}{p{13cm}  x{4.5cm}}
	\textbf{#2} - \textcolor{bgcol}{#3} &   \vspace{2.5pt}\textcolor{sectcol}{#1}
\end{tabular*}

\vspace{-8pt}
\textcolor{softcol}{\hrule}
\vspace{6pt}

	\foreach \desc in {#4}{
		$\cdot$ \desc\\[3pt]
	}
	
\eventspace
}


\newcommand{\cveventextra}[4]
{

\begin{tabular*}{1\textwidth}{p{13cm}  x{4.5cm}}
	\textbf{#2} - \textcolor{bgcol}{#3} &   \vspace{2.5pt}\textcolor{sectcol}{#1}
\end{tabular*}

\vspace{-4pt}
\textcolor{softcol}{\hrule}
\vspace{6pt}

	\foreach \desc in {#4}{
		$\cdot$ \desc\\[3pt]
	}
	
\eventspace
}

% creates a stretched box as 
\newcommand{\cveventmeta}[2]
{
	\mbox{\mystrut \hspace{87pt}\textit{#1}}\\
	#2
}

%----------------------------------------------------------------------------------------
% CUSTOM STRUT FOR EMPTY BOXES
%----------------------------------------- -----------------------------------------------
\newcommand{\mystrut}{\rule[-.3\baselineskip]{0pt}{\baselineskip}}

%----------------------------------------------------------------------------------------
% CUSTOM LOREM IPSUM
%----------------------------------------------------------------------------------------
\newcommand{\lorem}
{Lorem ipsum dolor sit amet, consectetur adipiscing elit. Donec a diam lectus.}



%============================================================================%
%
%
%
%	DOCUMENT CONTENT
%
%
%
%============================================================================%
\begin{document}


%use our custom fancy header definitions
\pagestyle{fancy}	





%----------------------------------------------------------------------------------------
%	HEADER IMAGE
%----------------------------------------------------------------------------------------

%\begin{figure}[H]
%\begin{flushright}
%	\includegraphics[trim= 320 130 460 210,clip,width=0.2\linewidth]{myfoto.jpg}	%trimming relative to image size!
%\end{flushright}
%\end{figure}


%---------------------------------------------------------------------------------------
%	TITLE HEADLINE
%----------------------------------------------------------------------------------------
\vspace{-8pt}
\begin{center}
	\HUGE \textsc{Paulo Fylippe Sell} \textcolor{sectcol}{\rule[-1mm]{1mm}{0.9cm}} \textsc{Resume}\\[2pt]
	
\end{center}



\vspace{6pt}


%---------------------------------------------------------------------------------------
%	META SECTION
%----------------------------------------------------------------------------------------
\metasection{Santa Catarina, Brazil}{\textbf{Status:} Telecommunications Engineer}
\metasection{paulosell.github.io}{\textbf{Fields:} Embedded Systems, Hardware Acceleration, Signal Processing}
\metasection{sellpaulof@gmail.com}{\textbf{Techs:} C/C++, Python, VHDL, Matlab, Java, Git}
\metasection{+55 48 984053102}{\textbf{Links:} \href{https://paulosell.github.io}{paulosell.github.io} $\cdot$ \href{https://github.com/paulosell}{github.com/paulosell}}
\vspace{-2pt}
\textcolor{softcol}{\hrule}
\vspace{6pt}

\normalsize

%---------------------------------------------------------------------------------------
%	SUMMARAY (optional)
%----------------------------------------------------------------------------------------
\vspace{-6pt}
\cvsection{Summary}
Telecommunications Engineer living in Santa Catarina, Brazil. Graduated from the Federal Institute of Santa Catarina with a thesis focused on the development of a framework for secure and reliable firmware updates for microcontroller-based embedded devices. Currently pursuing a master's degree in Applied Computer Science at the University of Vale do Itajaí (UNIVALI), specializing in Embedded Systems and conducting research with the objective of improving energy efficiency and increasing throughput in FPGA-based hardware accelerators for deep neural networks.\\

%============================================================================%
%
%	CV SECTIONS AND EVENTS (MAIN CONTENT)
%
%============================================================================%

%---------------------------------------------------------------------------------------
%	EXPERIENCE
%----------------------------------------------------------------------------------------

\cvsection{Experience}

\cvevent{feb 2022 - present}{Product Development Analyst}{Intelbras}{
	{Responsible for the backend development of embedded code in home \& office routers, including implementing new product features, managing device behaviors, and providing real-time device status information through REST APIs with HTTP and MQTT broker on Linux-based embedded solutions.},
	{Development with C and Shell script for hardware solutions delivered by suppliers based on SDK with embedded Linux and FreeRTOS.},
	{Experiences in development stack such as: TCP/IP, CMake, CMocka, Jansson.}
	{Linux driver and kernel development},
	{Development of ARM microcontroller-based products for Alexa Voice Service (AVS) solutions and the ALEXA/AWS ecosystem (IAM, S3, IoT Core).}
}

\cvevent{may 2021 - feb 2022}{Software Developer / Reseacher}{Federal Institute of Santa Catarina}{
	{Developing technical specifications and requirements for a Fiscal Authorization Device (DAF) for Electronic Consumer Invoice (NFC-e) usage in Santa Catarina.},
	{Tasks included description of architecture, communication, and secure hardware requirements; description of device use cases; writing of DAF's API and communication protocol; development of proof-of-concept firmware for DAF; integration of DAF with complementary software. Also involved in the creation of the set of instructions and rules for DAF certification.},
	{Development of proof-of-concept projects using C++, Java, and Python.},
	{Experiences with stacks such as: Flask, Docker, Redis, MySQL, Unity, USB Protocol, public key encryption, JWT.}
}
%
%---------------------------------------------------------------------------------------
%	EDUCATION SECTION
%--------------------------------------------------------------------------------------

\cvsection{Education}

\cvevent{2023 - present}{Master's degree in Applied Computer Science}{University of Vale do Itajaí }{
	{Specialization in Embedded Systems.},
	{Energy efficiency improvement and throughput enhancement in FPGA-based hardware accelerators for deep neural networks.}
}

\cveventextra{2015 - 2021}{Bachelor's degree in Telecommunications Engineering}{Federal Institute of Santa Catarina}{
	{Development of a framework for secure and reliable firmware update for microcontroller-based embedded devices.},
	{
		Award for \href{https://sol.sbc.org.br/index.php/sbseg_estendido/article/view/17354/17192}{best full paper} at the 15th Workshop on Scientific Initiation and Undergraduate Research (WTICG) of the 21st Brazilian Symposium on Information Security and Computer Systems (SBSeg 2021).
	}
}

% \cvevent{2011 - 2014}{Vocational High School}{Federal Institute of Santa Catarina}{
% 	{Telecommunications technician}
% }
\newpage
\cvsection{Academic Experience}

\cvevent{2020 - 2021}{Research Student}{Federal Institute of Santa Catarina}{
	{Development of the specification and the proof of concept of an embedded low-cost secure fiscal authorizing device for the state of Santa Catarina.}
}

\cvevent{2019 - 2019}{Research Student}{Federal Institute of Santa Catarina}{
	{Development of a DSP block for an Modified Discrete Cosine Transform (MDCT).},
	{Mostly used Matlab and VHDL}}



\cvsection{Fields of Knowledge}

$\cdot$  Embedded systems, wireless communication, protocol design, deep neural networks, RF, electric circuits, telecommunication systems, distributed systems.


%\textcolor{softcol}{\hrule}

%

%--------------------------------------------------------------------------------------------------
%	ARTIFICIAL FOOTER (fancy footer cannot exceed linewidth) 
%--------------------------------------------------------------------------------------------------



%============================================================================%
%
%
%
%	DOCUMENT END
%
%
%
%============================================================================%
\end{document}
